% Part: first-order-logic
% Chapter: beyond
% Section: introduction

\documentclass[../../../include/open-logic-section]{subfiles}

\begin{document}

\olfileid{fol}{byd}{int}

\olsection{مرور کلی}

منطق مرتبهٔ اول تنها دستگاه منطقی شایان توجه نیست: بسط‌ها و گونه‌های
فراوانی از منطق مرتبهٔ اول وجود دارد. یک منطق معمولاً از تعیین صوری یک
زبان، غالباً اما نه همیشه یک دستگاه استنتاج، و غالباً اما نه همیشه یک
معناشناسی مقصود تشکیل می‌شود. بااین‌همه، کاربرد فنی اصطلاح «منطق»
پرسشی آشکار پیش می‌کشد: منطق‌هایی که منطق مرتبهٔ اول نیستند چه نسبتی
با واژهٔ «منطق» در معنای شهودی یا فلسفی آن دارند؟ همهٔ دستگاه‌هایی که
در ادامه توصیف می‌شوند برای مدل‌سازی صورتی از استدلال طراحی شده‌اند؛
آیا می‌توانیم بگوییم چه چیزی آن‌ها را منطقی می‌سازد؟

پاسخ آسانی در کار نیست. واژهٔ «منطق» به شیوه‌های گوناگون و در بافت‌های
متفاوت به کار می‌رود و این مفهوم، همانند مفهوم «صدق»، از مواضع فلسفی
پرشماری تحلیل شده است. برای نمونه، ممکن است هدف استدلال منطقی را تعیین
این بدانیم که کدام گزاره‌ها بالضروره صادق‌اند، پیشینی صادق‌اند، مستقل
از تعبیر اصطلاحات غیرمنطقی صادق‌اند، به اقتضای صورت خود صادق‌اند، یا
به قرارداد زبانی صادق‌اند؛ و هر یک از این تلقی‌ها نیازمند توضیحی مفصل
است. حتی اگر توجه خود را به آن نوع منطقی محدود کنیم که در ریاضیات به
کار می‌رود، دربارهٔ دامنهٔ آن توافق چندانی وجود ندارد. برای نمونه، راسل
و وایتهد در \textit{اصول ریاضیات} کوشیدند، در سنت {em منطق‌گرایی} که
فرگه آغاز کرده بود، ریاضیات را بر پایهٔ منطق بپرورانند. دستگاه منطقی
آنان صورتی از منطق انواع بالاتر بود که به دستگاه توصیف‌شده در ادامه
شباهت داشت. آنان سرانجام ناچار شدند اصول موضوعی‌ای وارد کنند که بنا بر
بیشتر معیارها صرفاً منطقی به نظر نمی‌رسند (به‌ویژه اصل موضوع نامتناهی و
اصل موضوع تحویل‌پذیری)، اما بااین‌حال می‌توان بر این نظر بود که برخی
صورت‌های استدلال مرتبهٔ بالاتر را باید منطقی پذیرفت. در مقابل، کواین،
که هستی‌شناسی‌اش «قضایا» را اشیای مشروع گفت‌وگو نمی‌پذیرد، استدلال
می‌کند که منطق مرتبهٔ دوم و مراتب بالاتر در واقع نمودهایی از نظریهٔ
مجموعه‌ها در لباس مبدل‌اند؛ به بیان دیگر، دستگاه‌هایی که در آن‌ها روی
محمول‌ها کمیت‌گذاری می‌شود صرفاً منطقی نیستند.

فعلاً بهتر است چنین مسائل فلسفی‌ای را برای روزی دیگر بگذاریم و
دستگاه‌های زیر را صرفاً آرمانی‌سازی‌های صوریِ گونه‌های مختلف استدلال،
اعم از منطقی و غیرمنطقی، بدانیم.

\end{document}
